\section{Set Up}

Set up the map, Game Turn Marker, and train hexes normally. All counters begin
on their ordered sides unless stated otherwise.

\subsection{Royalist Army}

Set up \textbf{every counter of Rupert's and Newcastle's forces} in the set-up
hex printed on its counter. This includes Newcastle, Eythin, all seven Newcastle
Foot, all Northern Horse, and all four Royalist artillery pieces. The earlier
arrival has already happened; the Royalist army is assembled when play begins.

\subsection{Allied Rearguard}

Set up the following counters in their printed set-up hexes:

\begin{itemize}
  \item every Allied Heavy Horse, Light Horse, and Dragoon unit;
  \item Manchester, Cromwell, Leven, Fairfax, and Thomas Fairfax; and
  \item no Allied Foot or artillery.
\end{itemize}

Place all Allied Foot and all six Allied artillery pieces off-map, sorted into
the groups shown below.

\section{Allied Arrival Schedule}

Arriving pieces enter during their scheduled Allied March Phase through the
\textbf{southwest entry area}. This consists of south-edge hexes
0101--0108 and west-edge hexes 0101--0801.

The first map hex entered costs its normal terrain cost; there is no additional
cost for crossing the map edge. A piece need not enter on its scheduled turn,
but remains available on every later turn until it enters.

\begin{center}
\small
\begin{tabularx}{\textwidth}{@{}p{.72in}p{.72in}Y@{}}
\toprule
Turn & Time & Arriving pieces \\
\midrule
2 & 12:30 & \textbf{Fairfax Foot:} Rigby, Dodding, Ashton, Fairfax, Bright, and Constable. \\
3 & 1:00 & \textbf{Manchester Foot:} Manchester, 1/Crfrd, and 2/Crfrd. \\
4 & 1:30 & \textbf{Leven's first group:} Reserve, Sinclair, Erskine, Dunhope, Yester, Livgstn, Coupar, and Dunfer; \textbf{all six Allied artillery pieces}. \\
5 & 2:00 & \textbf{Leven's second group:} Killhead, Cassilis, Bchlch, Loudon, Rae, Hamilton, Maitland, and Fifeshire. \\
\bottomrule
\end{tabularx}
\end{center}

\begin{scenarioBox}{Approach Column}
On the turn it enters, an ordered Foot unit has 8 MPs instead of its printed
allowance. It pays normal terrain and hexside costs, and may move through
friendly pieces normally, but it may not enter an ordered Royalist unit's ZOC.
\newline
\newline
An entering Foot unit may not attack during that turn's Allied Combat Phase.
It is otherwise an ordered unit: it exerts a ZOC and defends at full strength.
\end{scenarioBox}

\begin{scenarioBox}{Arriving Artillery}
An arriving Allied artillery piece may move up to 8 MPs during its entry March
Phase, paying Foot movement costs. It may move only through clear or hilltop
(aka rough) hexes, or along lanes, and may not enter an ordered Royalist unit's
ZOC. Once this movement ends, the piece is unlimbered and cannot move again.
\newline
\newline
An arriving artillery piece cannot bombard until a later Allied Artillery
Phase. Thus the guns entering on Turn 4 can first bombard on Turn 5, at 2 p.m.
\end{scenarioBox}
